\section{Introduction}
\label{sec:introduction}

Network security remains a central challenge for modern computing systems. Traditional intrusion detection systems (IDS) rely heavily on signatures or static statistical rules, which limits their ability to cope with zero-day attacks, distribution shift, and adversarial manipulation \cite{liao2024survey, moustafa2023explainable, wang2023robustnessnids}. Recent learning-based IDS research has therefore moved toward self-supervised generalization \cite{golchin2024ssclids, xu2024gnnssl}, open-set recognition \cite{qiu2024vaemax}, and transformer-based feature modeling \cite{xi2024idsmtran}. At the same time, recent evidence shows that strong in-distribution benchmark results can degrade sharply under cross-dataset testing or stronger perturbations \cite{cantone2024crossdataset, wang2023robustnessnids}. Reinforcement learning (RL) offers a natural route toward adaptive defense, but many existing RL-based IDS pipelines still train against either static disturbances or simultaneously updated opponents, without explicitly approximating an attacker's best response.

This paper proposes a \textbf{Bi-level Adversarial Reinforcement Learning (Bi-ARL)} framework to address this gap and then extends the same bilevel idea to stronger supervised tabular detectors. Unlike standard multi-agent RL, where both agents are optimized simultaneously, our approach models the interaction as a \textbf{Stackelberg game}: the defender acts as the leader and is updated against an attacker that is repeatedly adapted toward a stronger response. The purpose is not to claim a solved robust IDS, but to test whether explicit leader-follower optimization yields a better detection trade-off than ordinary RL baselines and whether the same training principle remains useful after replacing the weak RL detector backbone with stronger supervised models.

\subsection{Contributions}
Our main contributions are:
\begin{enumerate}
    \item We formulate IDS training as a bi-level attacker-defender optimization problem and implement a practical PPO-based Bi-ARL pipeline for tabular network-flow benchmarks.
    \item We show on NSL-KDD that Bi-ARL achieves the best mean F1 score among the RL baselines while substantially reducing FPR relative to vanilla PPO.
    \item We refactor the original RL line into a route-C family of bilevel adversarially trained supervised detectors, including BiAT-MLP and BiAT-FTTransformer, in order to preserve the core bilevel idea while using stronger tabular backbones.
    \item We extend the evaluation to UNSW-NB15 and CIC-IDS2017, where BiAT-FTTransformer reaches 89.51\% mean F1 and 14.97\% mean FPR on UNSW-NB15, becoming the strongest neural detector in our study and achieving the lowest false positive rate among all compared methods on that dataset.
    \item We provide a deliberately conservative empirical analysis: the original RL formulation remains useful mainly on NSL-KDD, while the refactored route-C framework is materially stronger on modern data but still trails the strongest boosted-tree baselines on overall F1.
\end{enumerate}

